\documentclass[10pt,a4paper]{amsart}
\usepackage[margin=19mm]{geometry}
\usepackage{amsmath,amssymb,mathtools}
\usepackage[scr=rsfs,cal=euler]{mathalfa}
\usepackage{xfrac}
\usepackage[unicode,hidelinks,bookmarks=false]{hyperref}
\newcommand{\ie}{i.e.\ }
\newcommand{\cf}{cf.\ }
\newcommand{\resp}{respectively\ }
\newcommand{\Subsection}{\subsection}
\providecommand{\nobreakdash}{\nobreak\hbox{-}}
\setlength{\emergencystretch}{2em}
\multlinegap0pt
\usepackage{tikz}
\usepackage{bm}
\usepackage{amssymb}
\usepackage{enumerate}
\newtheorem{lemma}{Lemma}
\newcommand{\m}{\mathcal}
\title{A sharp product bound for non-trivial cross-intersecting families}
\author{Yang Huang}
\date{Source: arXiv:2606.23322v1 (CC BY 4.0)}
\begin{document}
\maketitle

\noindent\emph{Source and licence.} A sharp product bound for non-trivial cross-intersecting families. Version arXiv:2606.23322v1 (CC BY 4.0), distributed by arXiv under the Creative Commons Attribution 4.0 International licence, http://creativecommons.org/licenses/by/4.0/, as recorded in the arXiv licence metadata for this exact version and shown on the abstract page https://arxiv.org/abs/2606.23322v1. Full licence notice: \texttt{licenses/arxiv-2606.23322-CC-BY-4.0.txt}.\par\smallskip

\noindent\textit{Editorial scope.} Counted unit: the article's lemma end2 (source0.tex lines 2643--2660). The article's complete original proof of it is delivered with it (source0.tex lines 2662--2724, one \texttt{proof} environment of 63 lines). No mathematics is written, repaired or re-derived: the statement and the proof are the article's own lines, in the article's own notation, and no retained line is substituted or re-flowed.  No further source-local context is needed: every cross-reference of every retained line resolves inside the two delivered spans.  Nothing was omitted from any retained span, so the package carries no omission record.  No retained line contains a citation, so the wrapper carries no bibliography and no bibliography entry.  No retained line contains a figure environment, an image-inclusion command or drawing code, so no image is delivered and the package declares no asset dependency.  Editorial formatting only, in addition: an AMS \texttt{amsart} preamble that restates the article's own declarations and macros used by the retained lines, namely the environments \texttt{lemma} on separate counters, together with the standard packages those lines need.  The retained environments print in this document as Lemma 1 in order of appearance, whereas the article prints the same items with the numbering of its own full document.  The complete native source of the article as distributed in the arXiv e-print for this exact version is bundled unchanged as \texttt{sources/203422/source0.tex}.
\begin{lemma}\label{end2}
Let $n\ge 7$, $\m A, \m B \subset {[n]\choose 3}$ be non-trivial cross-intersecting families. 
If $|\m B(1)|\ge n-1$, then $|\m B(1)|+|\m A(\bar1)|\le 3n-8$, and equality holds only in one of the following three cases.
\begin{itemize}
\item[\rm (i)]
There exists $3$-set $S\subset [2,n]$ such that
$\m A(\bar1)=\{S\}$ and 
$\m B(1)=\left\{E\in\binom{[2,n]}{2}:E\cap S\ne\emptyset\right\}$.
\item[\rm (ii)]
There exists $2$-set $P\subset [2,n]$ such that $\m B(1)=\left\{E\in\binom{[2,n]}{2}:E\cap P\ne\emptyset\right\}$ and
$\m A(\bar1)=\{E\in \binom{[2,n]}{3}:P\subset E\}$.
\item[\rm (iii)]
There exist $x\in [2,n]$ and $P\in\binom{[2,n]\setminus\{x\}}2$ such that
$\m A(\bar1)=\{E\in \binom{[2,n]}{3}:x\in E, P\cap E\ne\emptyset\}$ and 
$\m B(1)=\left\{E\in\binom{[2,n]}{2}:x\in E\right\}\cup\{P\}$.
\end{itemize}
Moreover, $|\m A(\bar1)|\le 2n-7$. 
\end{lemma}

\begin{proof}
Note that
every set of $\m A(\bar 1)$ is a $3$-element vertex cover of
$\m B(1)$. Thus, $\tau(\m B(1))\le 3$.
Since $|\m B(1)|>n-2$, $\tau(\m B(1))\ge 2$. Together with $\tau(\m B(1))\le3$, we have $\tau(\m B(1))\in\{2,3\}.$ 

{\bf Case 1: $\tau(\m B(1))=2$. }
Let $\{u,v\}$ be a minimum vertex
cover. Put 
\[
a:=|\m B[\{1,u\}\overline{\{v\}}]|,
\quad
b:=|\m B[\{1,v\}\overline{\{u\}}]|,
\quad
\eta:=|\m A(\overline{\{1,u,v\}})|,
\quad
\varepsilon:=
\begin{cases}
1,& \text{ if } \{1,u,v\}\in \m B,\\
0,& \text{otherwise}.
\end{cases}
\]
Then $|\m B(1)|=a+b+\varepsilon,$ $a,b\ge1$, $|\m A(\{u,v\}\overline{\{1\}})|\le n-3$, $a,b \le n-3$,
and if $\eta=1$, then $\varepsilon=0$.

First, consider the case $a,b\ge3$. If $\varepsilon=1$, then every set in $\m A(\bar1)$ contains $\{u,v\}$, i.e. $|\m A(\bar1)|\le {n-3}$, so $|\m A(\bar1)|+|\m B(1)|\le {n-3}+a+b+1\le 3(n-3)+1=3n-8$.
If $\varepsilon=0$, then either $\eta=1$, in which case $a=b=3$, $|\m B(1)|=6$, and $|\m A(\bar1)|\le n-2$, and hence $|\m B(1)|+|\m A(\bar1)|<3n-8$; or $\eta=0$, in which case 
$|\m A(\bar1)|\le n-3$ and $|\m B(1)|\le 2(n-3)$, and hence, 
$|\m B(1)|+|\m A(\bar1)|\le 3(n-3)<3n-8$. Thus, if $a,b\ge3$, then $|\m B(1)|+|\m A(\bar1)|\le 3n-8$ with equality only if $a=b=n-3$ and $\varepsilon=1$, which is case \rm(ii). Moreover, $|\m A(\bar1)|\le n-2<2n-7$. 

Next consider the case $a=2$ or $b=2$. If, say, $b=2$, then $a\ge n-4$, $\eta=0$, $|\m A(\bar1)|\le n-3+1=n-2$, and $|\m B(1)|\le 1+2+n-3=n$, thus,
$|\m B(1)|+|\m A(\bar1)|\le n-2+n<3n-8$. In this case, $|\m A(\bar1)|\le n-2<2n-7$. 

Finally consider the case $a=1$ or $b=1$. If, say, $b=1$, then $|\m B(1)|\ge n-1$ and  $a\le n-3$ forces $a= n-3$ and $\varepsilon=1$. 
Consequently, $|\m B(1)|= n-1$ and $|\m A(\bar1)|\le 2n-7$,
and case \rm(iii) follows. 

{\bf Case 2: $\tau(\m B(1))=3$.}
Let $S=\{x_1,x_2,x_3\}$ 
be a minimum vertex cover. 
For each $i\in[3]$, let $a_i:=|\m B(\{1,x_i\}\overline{S\setminus \{x_i\}})|$, and let $q:=|\{B\in \m B(1): |B\cap S|=2\}|$.
Since $S$ is a minimum cover, $a_i\ge1$ for every $i$, and $|\m B(1)|=a_1+a_2+a_3+q.$
Besides $S$ itself, a three-element cover containing exactly two members of $S$ is obtained by omitting $x_i$ precisely when $a_i=1$. Let $s:=|\{i:a_i=1\}|$. Then $s=|\{A\in \m A(\bar1): |A\cap S|=2\}|$.
Let $t:=|\{A\in \m A(\bar1): |A\cap S|=1\}|$.
Then $t+q\le 3$.
Let $z:=|\{A\in \m A(\bar1): A\cap S=\emptyset\}|$. Then $z\in\{0,1\}$.
%A cover containing exactly one member of $S$ can occur at most once for each missing edge of $S$; denote the number of such covers by $t$. 
Thus $|\m A(\bar1)|=1+s+t+z$.
Let $r:=|[2,n]\setminus S|$. Note that $s\le\sum_{i=1}^3(r-a_i).$ 
Then $|\m A(\bar1)|=1+s+t+z \le 1+ \sum_{i=1}^3(r-a_i)+3-q.$
Indeed, this is immediate when $z=0$. If $z=1$, then $q=0$; the only additional exceptional possibility is $t=3$, in which case the three sets $a_1=a_2=a_3=1$, and the displayed inequality is again strict.
Therefore 
\[
|\m A(\bar1)|+|\m B(1)|\le 1+ \sum_{i=1}^3(r-a_i)+3-q+a_1+a_2+a_3+q=3r+4=3n-8.
\]
If equality holds, then $s=\sum_{i=1}^3(r-a_i).$ 
Since $r=n-4\ge3$, this forces $a_1=a_2=a_3=n-4.$
It then follows that $t=z=0$ and $q=3$. Thus $\m B(1)$ consists of all
edges meeting $S$, and $S$ is its unique three-element vertex cover.
This is case \rm(i).
Note that $|\m A(\bar1)|\le 1+ \sum_{i=1}^3(r-a_i)+3-q$ and $|\m B(1)|=a_1+a_2+a_3+q$, since $|\m B(1)|\ge n-1$, $|\m A(\bar1)|\le 3n-8-|\m B(1)|\le 2n-7$.
This completes the proof.
\end{proof}

\end{document}
